% document_id: BJP-THEORY
% revision: 0.59.0
% as_of: 2026-10-02
% [更新:0.59.0] [確認:0.59.0] — THEORY-CORE / REV-590-GUIDE-THEORY
% 0.58 scope: annual-only UTC quarters from raw samples; half/month quantiles and other physics retained, not fully revalidated.
% Editable source; PDF and search text are generated by tools/build_guides.py.
\documentclass[11pt,a4paper]{article}
\usepackage[margin=19mm,top=22mm,bottom=21mm,headheight=15pt]{geometry}
\usepackage{fontspec,xeCJK}
\setmainfont{Noto Sans CJK JP}
\setsansfont{Noto Sans CJK JP}
\setmonofont{Noto Sans Mono CJK JP}[Scale=0.86]
\setCJKmainfont{Noto Sans CJK JP}
\setCJKsansfont{Noto Sans CJK JP}
\setCJKmonofont{Noto Sans Mono CJK JP}[Scale=0.86]
\usepackage{amsmath,amssymb,longtable,booktabs,array,xcolor,xurl}
\definecolor{ink}{HTML}{183046}
\definecolor{accent}{HTML}{176D86}
\usepackage{fancyhdr,enumitem,fvextra}
\usepackage{tikz}\usetikzlibrary{arrows.meta,positioning,calc}
% Use standard article sectioning so page breaks and PDF anchors stay together.
\usepackage[unicode,colorlinks=true,linkcolor=accent,urlcolor=accent]{hyperref}
\usepackage{bookmark}
\DefineVerbatimEnvironment{verbatim}{Verbatim}{breaklines=true,fontsize=\small}
\newcommand{\code}[1]{\texttt{\detokenize{#1}}}
\newcommand{\GuideBuildID}{BJP-059-THEORY-DOC-20261002}
\renewcommand{\contentsname}{目次}
\setcounter{tocdepth}{2}
\setlength{\parindent}{0pt}
\setlength{\parskip}{5pt}
\setlength{\emergencystretch}{3em}
\linespread{1.08}
\renewcommand{\arraystretch}{1.25}
\setlist{nosep,leftmargin=1.8em}
\pagestyle{fancy}\fancyhf{}
\fancyhead[L]{\footnotesize Balloon-JP / 理論と数値計算}
\fancyhead[R]{\footnotesize 文書0.59 / 実装0.58}
\fancyfoot[L]{\scriptsize \GuideBuildID}
\fancyfoot[R]{\thepage}
\hypersetup{pdftitle={Balloon-JP 理論ガイド 文書0.59・実装0.58},pdfauthor={Balloon-JP},pdfsubject={Implemented model assumptions and finite-sample analysis}}
\begin{document}
\begin{titlepage}
\vspace*{18mm}{\color{accent}\rule{\linewidth}{1.2pt}}\par\vspace{10mm}
{\fontsize{27}{36}\selectfont\bfseries\color{ink}Balloon-JP 理論ガイド\par}
\vspace{8mm}{\Large 一飛行の式から、条件を比べる標本集合へ}\par
\vspace{12mm}文書改訂0.59 / 実装0.58　2026年10月2日\par
[更新:0.59.0] [確認:0.59.0]\par
{\small 文書改訂0.59は、固定0.58実装に対する資料監査で、地域平均への案内を訂正する。科学式・集計定義・実装は変更しない。0.58で第\ref{th:climate}節へ加えた原UTC標本の年間48区分、旧半月の境界、月の二つの中央値とp10--p90、格子重みとUTC選択の意味を継承する。細分化を独立標本数や気象の代表性の増加とは解釈しない。今回の読取は気象統計章と版入口の有限範囲であり、他の科学式を全面再検証した印ではない。数値・速度・実画面の受入と未確認はS36へ戻る。main統合は別判断である。}\par
\vfill
この本は、指定した場と仮定から一本の軌道を得る式と、明示した入力の仮幅から複数の飛行を比較する集計を説明する。
予報の正しさ、実機の係数、着地の運用精度は、その式を解けたこととは別に評価する。
0.52の追補は、第\ref{th:jra-b}小節以降に、原UTC別の仮想列と地上節点の接続を加えた。第\ref{th:climate}節の期間統計とは資料・目的・支持を区別する。
\par\vspace{5mm}
編集元：\path{docs/THEORY_GUIDE.tex}\par
対象コード：\path{balloon_sim/environment/}、\path{balloon_sim/flight/}、\path{balloon_sim/ensemble/}、\path{backend/climate/}\par
0.52の科学追補：\path{environment/model_surface.py}、\path{jra3q.py} と保存schema\par
文書改訂の比較元 0.58 C2：\path{95f09d2fc0ae076def05eb76df996b9f2c3b0ef4}\par
0.54の比較元 0.53 C2：\path{f2a3fa6fa0621076452001177ce3556ef189acd5}\par
0.52の比較元 0.51 C2：\path{f94027cec2ff56b732af7bd11b2df26bdd9513a5}\par
旧構造修復の比較元 C2-230：\path{c1623997d4d8142f9e34ebf496e4be5d9c442bf6}\par
ビルド識別子：\GuideBuildID
\end{titlepage}
{\small\setlength{\parskip}{0pt}\tableofcontents}\clearpage
% LOGIC-BEGIN:THEORY-CORE
\section{この計算は、何を入力すれば何を答えるか}\label{th:purpose}
日本の任意の地点と打上げ時刻を指定し、取得した気象場の中で、上昇、破裂、降下、地形との交差を一続きに求める。
北海道と和歌山の例は異なる場を同じ経路に通すための入力例であり、この二地点を対象範囲の上限にしない。
位置だけを変えられても、途中で気象の領域・時刻・高度を外れれば一本の飛行を支えられない。
そこで、\emph{軌道の式}と\emph{式が参照できる場の範囲}をともに計算の一部とする。

本版は、球面上で風と同じ水平速度を持つ気球を扱う。鉛直方向には、指定上昇率または内外等温の準定常浮力・抗力を選び、破裂後の降下則を別に選ぶ。
風の取得は計算前に終え、積分中の時刻・位置・高度から保存場を問い合わせる。この順序によって、同じ入力の再実行と、物理モデルだけを変える比較が可能になる。
GFS圧力面と地上量の接続は本版の実装近似である。JRAは0.51のstrict場を保持し、0.52では別schemaの暫定B方針で地上解析へ接続する（第\ref{th:jra-b}小節）。保存場はCLIと既存画面から同じ飛行核へ渡し、JRAの任意過去窓を選ぶ取得画面は未実装である。熱収支、鉛直慣性、Re依存抗力も未実装のまま区別する。

読む順序には依存がある。第\ref{th:coordinates}節で量の意味を固定し、第\ref{th:weather}節で任意の高度・時刻の気象値を作る。
その値を第\ref{th:vertical}節の鉛直式と第\ref{th:horizontal}節の水平式へ渡し、第\ref{th:solver}節で相の切替と停止を処理する。
最後に第\ref{th:validation}節で、各式を独立に確かめる道筋と現実に対して残る問いを示す。
第\ref{th:ensemble}節では、その一飛行を共有した元標本で繰り返し、経験範囲・干渉群・相別履歴へ結ぶ。単一の固定気象場に対する入力感度と、気象誤差を含む予報確率は別の問いである。
APIを探すときは構造ガイド、実行するときは\path{docs/COMMANDS.md}を読む。

\section{座標・時間・単位を先に固定する}\label{th:coordinates}
\subsection{積分する状態と、混同しやすい高さ}
状態は緯度$\phi$、経度$\lambda$、海面基準のモデル幾何高度$z$である。式では角度をradで書くが、APIとJSONはdegreeで受け渡しする。
時刻$t$は打上げからの秒数であり、気象照会時刻はUTCの打上げ時刻に$t$を足す。
UTC以外のオフセットやタイムゾーンなしの入力を黙って解釈しない。利用者がJSTで考えた時刻は入力作成時にUTCへ変換する。

\begin{longtable}{@{}p{.19\linewidth}p{.26\linewidth}p{.49\linewidth}@{}}
\toprule 量 & 単位・向き & 役割 \\\midrule\endhead
$u,v$ & m/s、東・北が正 & 水平移流速度。風向角のまま補間しない。\\
$z,H,z_g$ & m、gpm、m & 幾何高度、ジオポテンシャル高度、選んだ気象モデルの地形高度。\\
$p,T,q$ & Pa、K、kg/kg & 全圧、外気温、水蒸気質量/湿潤空気全質量。$q$は相対湿度ではない。\\
$w,V,A$ & m/s、m$^3$、m$^2$ & 鉛直速度、気球体積、投影面積。$w$は上向きを正。\\
$m_g,m_e,m_p$ & kg & ガス、膜、搭載物。上昇時総質量は三者の和。\\
$C_D,C_DA$ & 1、m$^2$ & 抗力係数、降下時の有効抗力面積。\\
$g,R,R_d,R_v$ & SI & 本版定数は順に9.80665、6371000、287.05、461.5。\\\bottomrule
\end{longtable}

気象のHGTをそのまま幾何的な距離とすると、高度に依存する誤差を持つ。本版は球対称重力$g(z)=g_0[R/(R+z)]^2$の積分から
\begin{equation}
 H=\frac{1}{g_0}\int_0^z g(\zeta)\,d\zeta=\frac{Rz}{R+z},\qquad
 z=\frac{RH}{R-H}
 \label{eq:height}
\end{equation}
を用いる。これは測地学的な局所ジオイド・緯度依存重力を完全に復元する式ではない。
30 kmで$z-H=z^2/(R+z)$は約141 mとなり、単位の違いを無視するより明示的である。
鉛直力式の$g$は一定なので、高さ変換と力計算を合わせた精密な重力モデルを主張しない。

APIの\code{height_agl_m}は$z-z_g(t,\phi,\lambda)$であり、測量された地表からの高さではない。
GFSの粗い地形、実地表DEM、海面、最低気圧面は別の面である。特に1000 hPaを地面と置かないことが、山地を含む日本の地点へ拡張するために必要になる。

\subsection{3D表示へ渡すときだけ、高度基準を換算する}\label{th:display-height}
KMLの\code{absolute}はEGM96ジオイド基準の高さ、Cesiumの位置指定はWGS84楕円体高である。
そこで、実結果のモデルASL高度$z$を描画するときだけ、同梱格子から得たジオイド高$N_{\mathrm{EGM96}}$ [m]を加える。
\begin{equation}
 h_{\mathrm{display}}=z+N_{\mathrm{EGM96}}(\lambda,\phi).
 \label{eq:display-height}
\end{equation}
$N>0$はジオイドが楕円体より上にあることを表す。43$^\circ$N, 141.5$^\circ$Eの格子点では$N=31.33$ mであり、$z=30000$ mは表示時に$h=30031.33$ mとなる。
これは球面近似のモデルASLを、表示上EGM96基準として扱う規約である。式\eqref{eq:height}を再適用したり、モデル高度を測量上の厳密な正標高へ改善したりする操作ではない。

\begin{center}\small
\begin{tabular}{@{}p{.24\linewidth}p{.33\linewidth}p{.37\linewidth}@{}}
\toprule 元の量 & 表示へ渡す高さ & 保持する区別 \\\midrule
実結果のモデルASL $z$ & $z+N_{\mathrm{EGM96}}$ & 入力・飛行結果・KMLの元の$z$は書き換えない。\\
人工模型の模式AGL $a$ & 読めた地物の楕円体高$h_g+a$ & 実飛行のASLやモデル地表からのAGLではない。\\
写真3Dから読めた地物高$h_g$ & 既に楕円体高 & $N$を再加算せず、モデル地表$z_g$とも置き換えない。\\\bottomrule
\end{tabular}
\end{center}
EGM96は全球15分格子の4隅を双一次補間する。これは地形の細部や実機の高度誤差を復元するものではない。
格子を読み込めなければ実ASLの立体表示を保留し、$N=0$へ黙って代替しない。原科学値はそのまま保持する。

地表投影と半透明カーテンは、空中の軌道と地上の位置の垂直な対応を読むための補助表示である。
実ASL経路では地物高度の取得を幕の前提にせず、全経路へ連続した幕を描く。KMLと共通の原点を保ち、通常は約1 km以下へ表示節点を加えて、同じ節点から空中線と幕の上端を作る。
表示節点数の上限を超える場合も、原節点の全経路を保持する。線と幕は同じ測地線補間・角度刻みを使い、疎な経路でも片方だけを直線弦としない。KML閲覧ソフト間の地形・補間・描画の差まで一致させるものではない。

幕の下端は、経路の最小表示高度を$h_{\min}$として、一定の楕円体高
\[
 h_{\mathrm{bottom}}=\min(-12000\,\mathrm{m},\ h_{\min}-1000\,\mathrm{m})
\]
まで延ばす。これは地形高度を推定する式ではなく、上端から下へ描くための表示上の下端である。
不透明な写真meshの奥にある部分は深度処理で隠れ、見える幕は軌道と地表投影の対応を示す。
meshが未読・欠落した場所では余長が見える場合があるが、そこを地面や衝突地点とは解釈しない。
下端には輪郭線を付けず、軌道全体へ視点を合わせる範囲にも地下の余長を含めない。
モデルの軌道が写真地物の下へ入っても原高度を持ち上げず、表示の重なりと数値積分の接地判定を分ける。
この表示のために科学結果、経験被覆や干渉群を再計算しない。

人工模型の模式AGLだけは、上端を決めるための地物高が必要である。
人工の全表示節点の高さが揃うまでは空中経路と幕を一括保留し、実ASLへ無言で読み替えない。
現在描画されたmeshからの有限照会なので、長い人工経路では全点が揃わない場合がある。
この制限を実結果の連続幕へ持ち込まず、人工背景の試験と実写真での確認も分けて記録する。
格子の来歴・利用条件は第10節の表示用資料へ戻る。

\subsection{風向表示は、成分から最後に作る}
気象風向は吹いてくる方位FROMで、北0度・時計回りである。$u>0,v=0$の東向き流れは西風270度となる。
必要なら$[\operatorname{atan2}(u,v)180/\pi+180]\bmod360$から表示する。
進行先TOは反対向きであり、359度と1度の算術平均を180度にしてはいけない。
本体は$u,v$を受け渡すため、この円環量の平均を計算経路へ持ち込まない。

\section{気圧面・モデル面から連続な気象場を作る}\label{th:weather}\label{THEORY-WEATHER}
最初の三小節はGFS圧力面と地上量の実装を説明する。JRAのstrict方式は第\ref{th:jra-model}小節、水平と鉛直の順を変える暫定B方式は第\ref{th:jra-b}小節で区別する。
\subsection{時刻を選ぶことと、値を補間することは別である}
GFSの初期時刻runに予報時間leadを加えたものがvalid時刻である。取得した時刻はそのどちらとも異なる。
本版は単一runを明示し、そのrunの0--120時間は1時間間隔、その後384時間までは3時間間隔という配布軸から、要求区間の両端を支える時刻を選ぶ。
正規の120/123時間の間は補間できるが、本来あるべき121時間等を勝手に作らず、通常配布軸の内部欠落を広い間隔の補間で隠さない。
要求区間が一つのvalidと完全一致するときだけ、その一点の照会を一枚で支えられる。

打上げをvalidに合わせる必要はない。必要なのは、全飛行中のすべての積分段に正の重みで寄与する気象標本があることである。
データの開始・終了や箱の境界は閉区間として扱うが、境界外は外挿しない。
ここで384時間までというのは選択可能な予報時間の仕様であり、同じ精度や現在の配布ファイルの存在を保証するものではない。

\subsection{先に各列を同じ高度へ再構成する理由}
気圧面の高度$H_k$は時刻と水平格子点によって違う。二つの時刻の$H_k$と風を先に平均してから鉛直補間する方法は、各時刻の同じ高度の風を平均する方法と一般に一致しない。
軌道が必要とするのは「今の幾何高度の場」なので、GFSとJRA strictは\emph{各時刻・各水平列で同じ高度へ再構成し、その値を水平・時間で混合する}。

\begin{center}
\begin{tikzpicture}[x=1cm,y=1cm,>=Latex,font=\small]
\path[use as bounding box](0,-.7) rectangle(15.5,4.6);
\fill[accent!5](.1,.1) rectangle(7.3,4.4);
\fill[accent!5](8.1,.1) rectangle(15.3,4.4);
\node[anchor=west,text=ink] at(.4,4.1){\textbf{時刻$t_0$の列}};
\node[anchor=west,text=ink] at(8.4,4.1){\textbf{時刻$t_1$の列}};
\draw[->,accent](1,.6)--(1,3.8) node[above]{$H$};
\draw[->,accent](9,.6)--(9,3.8) node[above]{$H$};
\draw[thick,accent](1.3,1.1)--(3.2,1.1) node[right]{$H_k(t_0),f_k(t_0)$};
\draw[thick,accent](1.3,3.2)--(3.2,3.2) node[right]{$H_{k+1}(t_0),f_{k+1}(t_0)$};
\draw[thick,accent](9.3,1.7)--(11.2,1.7) node[right]{$H_k(t_1),f_k(t_1)$};
\draw[thick,accent](9.3,3.6)--(11.2,3.6) node[right]{$H_{k+1}(t_1),f_{k+1}(t_1)$};
\draw[dashed,very thick,orange!85!black](.4,2.4)--(15,2.4);
\node[fill=white,text=orange!70!black] at(7.7,2.4){共通の$H(z)$};
\fill[orange!85!black](2.2,2.4) circle(2pt);
\fill[orange!85!black](10.2,2.4) circle(2pt);
\draw[->,orange!85!black](2.2,2.3)--(2.2,.45);
\draw[->,orange!85!black](10.2,2.3)--(10.2,.45);
\node[anchor=west,text=ink] at(2.5,.55){$f_{t_0}[H(z)]$};
\node[anchor=west,text=ink] at(10.5,.55){$f_{t_1}[H(z)]$};
\node[anchor=west,text=ink] at(.4,-.4){\textbf{次に混合：}各列の同じ高さの値に、水平位置と時間の重みを掛ける。};
\end{tikzpicture}
\end{center}
{\small\color{ink}図は時間方向だけを示す。実際には各時刻の周囲の水平格子列にも同じ処理を行う。気圧面の高さ自体を先に平均しない。}

上下一対の支持高度$H_k\le H\le H_{k+1}$に対し、
\begin{equation}
 \alpha=\frac{H-H_k}{H_{k+1}-H_k},\quad
 f_c(H)=(1-\alpha)f_{c,k}+\alpha f_{c,k+1},\quad
 p_c(H)=\exp[(1-\alpha)\ln p_k+\alpha\ln p_{k+1}]
 \label{eq:column}
\end{equation}
とする。$f$は$u,v,T,q$、添字$c$は時刻と水平位置を固定した列を表す。
圧力を対数補間するのは正値性と大気の概ね指数的な鉛直減少に合わせるためであり、各列の厳密な静水圧積分を行っているわけではない。

水平二方向と時間の線形重みの積を$\omega_c$とすると、最終値は
\begin{equation}
 f(t,\phi,\lambda,z)=\sum_c\omega_c f_c[H(z)],\qquad
 p(t,\phi,\lambda,z)=\sum_c\omega_c p_c[H(z)],\quad
 \sum_c\omega_c=1.
 \label{eq:mix}
\end{equation}
圧力も列で指数へ戻した値を混合するため、最後まで対数値を混ぜる方式とは異なる。
高度順序の逆転、正の重みを持つ支持値の欠測、上端不足は明示エラーになる。正確な格子/時刻端点では重みゼロの隣接点に依存しない。
これにより「偶然不要な隣の欠測」と「必要な支持の欠測」を分けられる。

\subsection{地表から最下層までの接続と、その代償}
各列で地表圧より高圧の面、地形より10 gpm以内の面を除き、最初の地上気圧面を選ぶ。
その面までを、地表圧、地上2 mの$T,q$、地上10 mの$u,v$に接続する。
圧力は地表から対数補間し、他の量は各アンカーから線形補間する。アンカーより下はその値を保持する。
風の対数則、粗度長、安定度補正、局地境界層を解いていないため、この接続は低高度の近似を明示して飛行をつなぐ選択である。

異なる標高の格子点を混合すると、照会点は混合地形より上でも、隣の高い列の地表より下になり得る。
本版はその\emph{寄与列だけ}を地表値へ延長し、\code{terrain_stencil_surface_clamp}と最大延長深さ\code{terrain_stencil_clamp_depth_m}を返す。
照会高度自体が混合地形より下なら\code{BELOW_GROUND}で拒否する。この規則は、地下の気圧面値を空中の風と解釈せずに地形境界を連続につなぐ近似である。
高低差が大きい場所で精度が良くなる根拠にはならない。

地上接続幅は\code{surface_bridge_height_m}、近似の種類は\code{quality}に返る。
接地用の地形はGFS地形を同じ水平・時間重みで補間した面であり、細密DEMによる着地ではない。
将来DEMを差し替える際にも、気象モデルの地表との標高差と支持不足は残るので、単に地形配列だけを置換すれば問題が消えるとはしない。


\clearpage
\subsection{JRAモデル面では、圧力を列ごとに作る}\label{th:jra-model}
同じモデル面番号でも、圧力は地表圧によって変わる。「保存された過去の場を、任意の放球時刻と高度で使えるか」を判断するには、面の番号、圧力、高さを分けなければならない。
0.51の初期接続は、保存済みJRA-3Qモデル面の100層・2時刻・2$\times$2格子を使う。これは将来の領域・時刻数の上限ではなく、任意日時を取得する機能とは別の有限な接続である。

地表側からモデル層を$k=1,\ldots,100$とする。半層（層境界）の係数$A$ [Pa]と$B$ [無次元]から、各時刻・水平列$c$の地表圧$p_{s,c}$を使って
\begin{equation}
 p_{c,k+1/2}=A_{k+1/2}+B_{k+1/2}p_{s,c},\qquad
 p_{c,1/2}=p_{s,c},\quad p_{c,100+1/2}=0
 \label{eq:jra-half}
\end{equation}
を作る。面の代表圧力（full level）は半層圧力の単純な算術平均にしない。
$L=p_{c,k-1/2}$、$U=p_{c,k+1/2}$、基準単位$p_*=1\,\mathrm{Pa}$とおくと、気象庁TL479仕様のSimmons--Burridge式（$C=1$）は
\begin{equation}
 \begin{aligned}
 p_{c,k}&=p_*\exp\!\left[
 \frac{L\ln(L/p_*)-U\ln(U/p_*)}{L-U}-1\right]
 &&(k=1,\ldots,99),\\
 p_{c,100}&=\frac{p_{c,99+1/2}}{2}.
 \end{aligned}
 \label{eq:jra-full}
\end{equation}
最上層は別定義なので、$U=0$を対数式へ入れない。式\eqref{eq:jra-full}は圧力座標の定義であり、気球の鉛直運動式ではない。
コードでは$U>0$に対して$r=(L-U)/U$とおき、等価な$L\exp[\log(1+r)/r-1]$を\code{log1p}で評価する。
係数は保存原本の既知版に固定し、層数だけ同じ未知の係数を別製品として無言で採らない。

\subsection{strict方式：各列のH上で補間し、任意UTCへ混合する}\label{th:jra-strict}
各モデル面には保存されたジオポテンシャル高度$H_{c,k}$がある。照会の幾何高度$z$を式\eqref{eq:height}で$H(z)$へ変換し、各正重み列の上下の面を選ぶ。
\begin{equation}
 \alpha_c=\frac{H-H_{c,k}}{H_{c,k+1}-H_{c,k}},\qquad
 f_c(H)=(1-\alpha_c)f_{c,k}+\alpha_c f_{c,k+1},
 \quad f\in\{u,v,T,q,p\}.
 \label{eq:jra-column}
\end{equation}
その後に式\eqref{eq:mix}と同じ水平・時間の重みで混合する。端点では重みゼロの隣接点を要求しない。
このJRA strict接続は\textbf{圧力もHに対して線形}に補間し、既存JRA試作の意味を保つ。GFSの式\eqref{eq:column}の対数圧力補間とは違う。
密度は得られた$p,T,q$から式\eqref{eq:density}で作り、密度の直接補間や、厳密な静水圧再積分を行ったとはしない。

\begin{center}\small
\begin{tabular}{@{}p{30mm}p{59mm}p{67mm}@{}}
\toprule 判断する箇所 & GFS圧力面 & JRAモデル面のstrict接続 \\\midrule
圧力座標 & 保存された気圧面の$p_k$ & 列の$p_s$と半層係数からfull圧力を生成 \\
列内の圧力 & $\log p$をHで線形補間 & $p$をHで線形補間 \\
地表付近 & 地上量の橋渡し・列の有限延長 & 各列の実支持を要求。橋渡し・延長なし \\\bottomrule
\end{tabular}
\end{center}
補間方法の差を、過去場へ接続したついでに変更しない。別方式を採る場合は、保存原値をそろえた比較と採否を別に示す。

\clearpage
\subsection{時刻窓と地表付近の支持が、飛行を制限する}\label{th:jra-support}
JRA解析のnative時刻は00・06・12・18 UTCである。打上げをこの刻みに丸める必要はなく、保存時刻が照会を挟めば02:17:13 UTCも内挿できる。
ただし「保存軸の隣接間隔が6時間」と「要求した飛行窓の先頭・末尾まで資料がある」は異なる。
例えば02:17:13から最大4時間を支えるには、00・06 UTCの2枚だけでは足りず12 UTCも必要になる。
現在の形式検査は保存軸の内部連続を調べる。要求窓全体の取得・事前被覆を保証する機能ではない。単一解析時刻を保存した場合は、その瞬間のみが時間支持になる。

地表のHを内挿して地形高度$z_g$を返せても、同じ高さの気象値が得られるとは限らない。
正の水平・時間重みを持つ列集合を$C_+$とし、各列の地表を$H_{g,c}$、最下モデル面を$H_{c,1}$とする。
欠測がなく、各面が地表より上の通常の列では、共通の鉛直支持下端は
\begin{equation}
 H_{\min}=\max_{c\in C_+}H_{c,1},\qquad
 \overline H_g=\sum_{c\in C_+}\omega_c H_{g,c},\qquad
 H\ge H_{\min}
 \label{eq:jra-support}
\end{equation}
を必要とする。これは全量の支持を保証する十分条件ではなく、各要求値の欠測・上端・時刻・水平範囲も別に確認する。
格子間では高い地形の列にも正の重みがあるため、$H_{\min}-\overline H_g$は一列の地表--最下層間隔より大きくなり得る。

\begin{center}
\begin{tikzpicture}[x=1cm,y=1cm,>=Latex,font=\small]
\path[use as bounding box](0,-.45) rectangle(15.8,3.65);
\fill[accent!5](.2,.1) rectangle(15.5,3.5);
\node[anchor=west,text=ink]at(.5,3.2){\textbf{低地列}};
\node[anchor=west,text=ink]at(5.9,3.2){\textbf{高地列}};
\draw[thick,ink](.6,.6)--(3.7,.6) node[right]{$H_{g,1}$};
\draw[thick,accent](.6,1.0)--(3.7,1.0) node[right]{$H_{1,1}$};
\draw[thick,ink](6.0,2.1)--(9.0,2.1) node[right]{$H_{g,2}$};
\draw[thick,accent](6.0,2.5)--(9.0,2.5);
\draw[dashed,thick,orange!85!black](.6,2.5)--(14.8,2.5);
\node[fill=accent!5,anchor=west,text=orange!70!black]at(10.0,2.7){共通下端 $H_{\min}=H_{2,1}$};
\draw[densely dotted,ink](.6,1.4)--(14.8,1.4);
\node[fill=accent!5,anchor=west]at(10.0,1.15){内挿地形 $\overline H_g$};
\draw[<->,thick,orange!85!black](10.1,1.43)--(10.1,2.47);
\node[anchor=west,text=ink]at(.5,-.25){模式図：内挿地形より上でも、高地列の気象がまだ支えない高さがある。};
\end{tikzpicture}
\end{center}

固定2時刻・2$\times$2格子では各列の地表--最下層差は約8 gpmでも、中央の水平・時間内部点で共通下端と内挿地形の差は約123 gpmとなる。
診断飛行の下端停止で得た約137 m AGLは、\emph{移流後の別の位置・時刻}における最後の有効記録の$z-z_g$である。123 gpmと137 mを同じ点の差や補間誤差として引き算しない。

旧strict選択肢は、列の地下・最下層未満を成功値へ補完しない。必要な列の欠測を捨てて重みを再正規化せず、支持外では停止する。
第\ref{th:solver}節の接地イベント探索でも、未支持の気象を作ることはできない。地表を$H_{\min}$へ引き上げたり、最後の有効点を着地点に置き換えたりしない。
この不足に対し、次小節の暫定Bは別の再構成規則を明示する。strictで停止した結果を、着地した結果へ書き換えるものではない。

幾何30 kmは式\eqref{eq:height}で約29859.397 gpmであり、30000 gpmと同じ高さではない。
一地点の幾何30 kmで全5量を照会できても、そこへ至る全軌道、時刻窓、水平移流先、地表付近を支えた証拠にはならない。
0.51で確認したのは固定原本の支持内再構成と停止の接続である。0.52の有限な完飛行と区別し、実機精度・過去数年の落下分散の受入へ広げない。個別窓の実行数値・停止履歴・比較はS36に置く。


\clearpage
\subsection{暫定B：原UTCごとに同じ面を水平合成する}\label{th:jra-b}
地上からの飛行を実過去場で比較するため、\code{native_horizontal_first_surface_v1}を明示的な暫定方針として加える。
旧strict場の保存形式・式\eqref{eq:jra-column}は保つ。Bは地表だけの修復ではなく、\textbf{水平と鉛直の順序を変え、上空の値にも影響する}。列別の地表延長を使うA案は比較実験に留め、本版の保存場の選択肢には含めない。

原UTCを$r$、水平格子列を$c$、native面を$k$とする。同一面のHと量$f$を、同じ非負の水平重み$w_c$で
\begin{equation}
 \bar H_{r,k}=\sum_c w_c H_{r,c,k},\qquad
 \bar f_{r,k}=\sum_c w_c f_{r,c,k},\qquad \sum_c w_c=1
 \label{eq:jra-b-horizontal}
\end{equation}
と合成する。nativeは保存された製品の同一モデル面番号を指す。NCARのregular Gaussian格子は配布側で水平補間された製品であり、無加工のJMA水平格子を意味しない。
地表H・SPと2m/10m解析量にも同じ水平重みを使う。

\begin{center}
\begin{tikzpicture}[x=1cm,y=1cm,>=Latex,font=\small]
\path[use as bounding box](0,.5) rectangle(15.6,3);
\node[draw=accent,fill=accent!5,rounded corners,text width=4.0cm,minimum height=2.1cm,align=left,anchor=north west,inner sep=7pt](h) at(.1,3){\textbf{1 原UTCの仮想列}\\native H/Fを水平合成\\地表と2m/10mも合成};
\node[draw=accent,fill=accent!5,rounded corners,text width=4.0cm,minimum height=2.1cm,align=left,anchor=north west,inner sep=7pt](v) at(5.5,3){\textbf{2 同じ絶対Hへ照会}\\仮想列の鉛直線形補間\\下層は変数別の地上橋};
\node[draw=accent,fill=accent!5,rounded corners,text width=3.9cm,minimum height=2.1cm,align=left,anchor=north west,inner sep=7pt](t) at(10.8,3){\textbf{3 最後に時刻を混合}\\原UTCごとの照会値を\\時間重みで線形補間};
\draw[->,thick,accent](h.east)--(v.west);\draw[->,thick,accent](v.east)--(t.west);
\end{tikzpicture}
\end{center}
{\small\textbf{変えないこと：}照会する絶対Hは全原UTCで共通。時刻間で面のHを先に平均してから鉛直照会する方式にはしない。}

仮想列と地上橋の鉛直照会を$\mathcal V_r(H)$、原UTCの時間重みを$\tau_r$とすると、
\begin{equation}
 f_B(t,\phi,\lambda,z)=\sum_r\tau_r\mathcal V_r[H(z)],\qquad \sum_r\tau_r=1.
 \label{eq:jra-b-time}
\end{equation}
各原UTCの仮想列がそのHを支える必要がある。時間平均した地形・上端だけを満たせばよい、という規則ではない。

圧力は\textbf{B1：元の各列で式\eqref{eq:jra-full}を評価してからfull $p$を水平平均}する。
平均SPからfull $p$を再計算する別方式は、非線形のため一般には同じでない。
鉛直は$p$自体をHで線形補間し、GFSの対数圧力規則へ切り替えない。

\clearpage
\subsection{地上節点と接続面を、量ごとに固定する}\label{th:jra-b-bridge}
本実装では地上解析のT/qを幾何AGL 2m、u/vを10m、SPを地表に置く。地表の直接観測や最下native面の別名とはしない。
原UTCの仮想地表$\bar H_{g,r}$から$z_{g,r}=R\bar H_{g,r}/(R-\bar H_{g,r})$を求め、AGL $a$の節点を
\begin{equation}
 H_{a,r}=\frac{R(z_{g,r}+a)}{R+z_{g,r}+a}
 \label{eq:jra-anchor}
\end{equation}
へ置く。単に$\bar H_{g,r}$へ2/10 gpmを加える操作ではない。

\begin{center}\small
\begin{tabular}{@{}p{19mm}p{31mm}p{28mm}p{69mm}@{}}
\toprule 量 & 低い節点 & 固定接続面 & その間・その下の意味 \\\midrule
$p$ & 地表SP ($a=0$) & level1 & 地表からnative1までH線形 \\
$T,q$ & 2m解析 ($a=2$) & level1 & 地表--2mは定値、2m--native1はH線形 \\
$u,v$ & 10m解析 ($a=10$) & level2 & 地表--10mは定値、10m--native2はH線形 \\\bottomrule
\end{tabular}
\end{center}
量$f$の接続面を$j_f$、低い節点の値を$D_{r,f}$とおく。地表以上・接続面未満では
\begin{equation}
 \mathcal V_{r,f}(H)=(1-\beta_f)D_{r,f}+\beta_f\bar f_{r,j_f},\qquad
 \beta_f=\max\!\left(0,\frac{H-H_{a_f,r}}{\bar H_{r,j_f}-H_{a_f,r}}\right).
 \label{eq:jra-bridge}
\end{equation}
接続面以上は、その面から上のnative点間で線形補間する。
節点以下の定値区間でも、固定した橋の上端値が存在することを要求する。上端が欠測の橋を、下の定値だけで成立したと扱わない。

\textbf{固定する対象は面IDである。} 構築時に、保存窓の全原UTC・全列で地形と各接続面のHが存在し、接続面が該当AGL節点より上であることを調べる。
照会ごとに「最初の10m超の面」を探し直さない。v1のID配置が成立しない窓は拒否し、別面へ自動変更しない。
水平合成後の節点順序も照会時に再検査する。

この風の再構成では\textbf{native level1のu/vを使わず}、10m解析とlevel2への橋で置き換える。
元配列は保持するが、その風の元節点値を再現する約束はしない。level1のp/T/qまで一緒に捨てない。
全列でlevel1が約8mとなった実窓は一例であり、8mを普遍的な層高としない。
地上解析を保持するこの定値/直線規則は、粗度・安定度・風の対数則を解く接地境界層モデルではない。
地上量の高さはJMA TL479仕様第4.1節、モデル面は第4.5節、単位/GPは第5.1節と保存DASへ照合した。製品の同定と、この接続規則の採否は別である。

\clearpage
\subsection{Bの支持・品質と、比較からは言えないこと}\label{th:jra-b-limits}
Bは各原UTCの仮想地表から仮想上端までを使う。上端は$\bar H_{r,100}$であり、strict方式の「寄与する元列の上端の最小値」と同じではない。
元の高地列の地表より低い高さを、仮想列で照会できる場合もある。これは\textbf{新しい場の定義}であり、strictの支持検査を通ったという説明へ替えない。

\begin{longtable}{@{}p{36mm}p{116mm}@{}}
\toprule 境界 & 保持する判定 \\\midrule\endhead
時刻・水平範囲 & 正の重みで必要な元UTC・格子が支持内にある。要求窓全体の事前被覆は保存軸の内部6h連続検査と別。 \\
地表・上端 & 全体の内挿地形と、各原UTCの仮想地表を下回れば拒否。仮想上端より上へ外挿しない。地下許容$10^{-7}$ gpmは高度変換の丸めだけで、地形列の有限延長ではない。 \\
固定joinの幾何 & 全窓の地形/接続高度欠測・節点順序不成立は構築時に拒否。補間で欠けたjoinを推定しない。 \\
照会に必要な値 & 正重みの必要値を平均前にも検査する。必要なnative H列の欠測、選択native圧力の欠測/地下判定、要求量の欠測は停止。欠けた列を除いて重みを再正規化しない。 \\
不要な値との区別 & 読込時の既知native量の妥当性検査は維持。地上診断値はその橋で要求する量を照会時に検査し、使わない地上量の欠測を必要値と混同しない。 \\\bottomrule
\end{longtable}

保存bundleは\path{balloon.weather.jra3q_surface/1}と明示policyを要求する。
旧\path{balloon.weather.jra3q_model/1}を自動変換しない。metadataにはpolicy・join ID・不使用のnative風・再構成順序を残す。
照会のqualityには方針を常設し、橋を使った場合は\code{surface_layer_bridge}、定値保持と最低風の不使用も区別する。
\code{surface_bridge_height_m}は照会に使った橋の最大幾何AGLで、誤差推定値ではない。

H/p/T/qの独立した合成・補間は、静力学整合や質量・エネルギー保存を保証しない。
地形は気象モデルの地形であり、細密DEM・建物・樹木に対する接触精度へ受入を広げない。
A/Bの着地点差は再構成に対する感度、刻み変更の差は有限な数値比較であり、実測飛行に対する精度順位や確率較正ではない。
実窓の完飛行・支持外停止・費用はS36へ記録し、本節の規則と区別する。

\clearpage
\section{鉛直モデルを組み合わせる}\label{th:vertical}\label{THEORY-DYNAMICS}
\subsection{初版で選べる経路と、まだ選べない経路}
上昇法則、破裂条件、降下法則は異なる問いへの答えである。したがって、一つのモデル名にすべてを固定せず、入力設定で組み合わせる。
\begin{longtable}{@{}p{.23\linewidth}p{.33\linewidth}p{.38\linewidth}@{}}
\toprule 選択箇所 & 実装済みの選択 & 必要量・条件 \\\midrule\endhead
上昇 & \code{constant_speed} & 指定した正の上昇率。風のみで上昇可能。\\
上昇 & \code{isothermal_buoyancy} & $p,T,q$、ガス・膜・搭載質量、ガス定数、一定$C_D$。\\
破裂 & \code{altitude} & 幾何高度しきい値。どちらの上昇にも使える。\\
破裂 & \code{diameter} & 直径しきい値。体積を持つ等温上昇に限る。\\
降下 & \code{constant_cda} & 破裂後有効質量、$C_DA$、気象密度。\\
降下 & \code{rated_speed} & 基準降下率と、気象密度または指数密度。\\\bottomrule
\end{longtable}
指定上昇・指定破裂高度・指数密度の基準降下率を組み合わせると、Tawhiri型の簡単な経路を残せる。
本版固有の取得、補間、半径、数値積分を使うため、CUSF固定原版との全出力同一性を意味しない。
Re依存抗力、慣性を持つ鉛直運動、熱状態、ガス損失、過圧の弾性モデルは未実装であり、未知のmodeを指定した場合は別モデルへ黙って代替しない。

\subsection{比湿から湿潤空気密度を作る}
乾燥空気質量$(1-q)m$と水蒸気質量$q m$が同じ$T,V$を持つ理想混合気体とする。
二つの分圧の和は$pV=m[(1-q)R_d+qR_v]T$なので、
\begin{equation}
 \rho_a=\frac{p}{T[(1-q)R_d+qR_v]},\qquad 0\le q<1
 \label{eq:density}
\end{equation}
となる。$q=0$では乾燥空気へ戻り、同じ$p,T$なら水蒸気の増加で密度が小さくなる。
凝結水・氷の質量を含む雲の混合物までは表さない。温度や圧力の欠測を標準大気で埋める処理もない。

\subsection{内外等温の浮力と抗力から上昇率を得る}
ガス質量$m_g$を一定、内部温度を外気温$T$、内部圧を外気圧$p$と置く。
膜の過圧と熱遅れを省いた理想気体から、
\begin{equation}
 V=\frac{m_gR_gT}{p},\qquad d=\left(\frac{6V}{\pi}\right)^{1/3},\qquad A=\frac{\pi d^2}{4}
 \label{eq:gas}
\end{equation}
を得る。形は球で、上昇中の全質量は$m=m_g+m_e+m_p$である。
浮力$\rho_a Vg$から重量$mg$を引いた自由揚力と、二乗抗力$\rho_a C_DAw^2/2$の釣合いを解くと、
\begin{equation}
 F_{\rm free}=g(\rho_a V-m),\qquad
 w_\uparrow=\sqrt{\frac{2g(\rho_a V-m)}{\rho_a C_DA}}
 \label{eq:ascent}
\end{equation}
となる。これは加速度が無視できる準定常近似であり、放球直後から終端速度にあると仮定する。
自由揚力が非正なら\code{INSUFFICIENT_LIFT}で停止する。ゼロを平方根へ入れて滞空を続けたり、指定速度へ切り替えたりしない。

$R_g$の既定2077.1 J/(kg K)はヘリウム相当のモデル値であり、純度や充填量の測定を代替しない。
質量と$C_D$は利用者入力で、例の数値は実機同定値ではない。$\rho_a V$が同じ温度・圧力の消去によって単純化しても、体積と面積は高度で変わるので上昇率まで一定とは限らない。
式を分けた理由は、将来熱収支や過圧を加えるとき、体積を返す関係式と必要状態量を差し替え、その影響を同じ浮力・抗力の収支で調べられるようにするためである。

\subsection{破裂は上昇と降下をつなぐ条件である}
高度破裂は$z-z_b=0$、直径破裂は$d-d_b=0$をイベント関数とする。
後者は式\eqref{eq:gas}の体積が必要なため、指定上昇率だけのモデルとは組み合わせない。
破裂で位置と時刻は連続に接続するが、鉛直速度は降下モデルの終端速度へ直ちに変わる。
破裂衝撃、パラシュート展開時間、残存膜の飛散を解くものではない。この不連続は入力を失った失敗ではなく、準定常な二相モデルの仮定である。

\subsection{降下率は、何を既知とするかで指定する}
破裂後の有効質量$m_d$と有効抗力面積$C_DA$が与えられる場合、浮力を省いた重量と抗力の釣合いは
\begin{equation}
 w_\downarrow=-\sqrt{\frac{2m_dg}{\rho_a C_DA}}.
 \label{eq:descentcda}
\end{equation}
ここでの$m_d$は残存ガス・膜の扱いを含めて利用者が定める。上昇時質量を自動で引き継ぐことはしない。

代わりに密度$\rho_0$での基準速度$v_0>0$を与えれば、同じ$C_DA/m_d$が維持される仮定のもとで
\begin{equation}
 w_\downarrow=-v_0\sqrt{\frac{\rho_0}{\rho_a}},\qquad
 C_DA=\frac{2m_dg}{\rho_0v_0^2}
 \label{eq:rated}
\end{equation}
と書ける。この二指定は物理的に重複する自由度なので、同時に独立パラメータとして掛け合わせない。
入力のしやすさと、どの量に根拠があるかで選ぶ。

簡易経路では気象密度の代わりに$\rho(z)=\rho_0\exp(-z/H_s)$を採用し、
\begin{equation}
 w_\downarrow=-v_0\exp\left(\frac{z}{2H_s}\right),\qquad
 t_d=\frac{2H_s}{v_0}\left[\exp\left(-\frac{z_g}{2H_s}\right)-\exp\left(-\frac{z_b}{2H_s}\right)\right]
 \label{eq:exponential}
\end{equation}
となる。右の着地時間は地面が一定の場合の解析解であり、独立試験にも使う。
$H_s$既定8400 mは経験的な簡易値で、実際の温度プロファイルを表す値ではない。
気象密度との比較では同じ風・破裂高度を保つと、降下法則の効果を分離できる。

\section{水平移流を球面上で積分する}\label{th:horizontal}
水平慣性を省き、気球はその場の東西風と南北風に乗る。緯度方向の弧長は$(R+z)d\phi$、経度方向は$(R+z)\cos\phi\,d\lambda$なので、
\begin{equation}
 \dot\phi=\frac{v(t,\phi,\lambda,z)}{R+z},\qquad
 \dot\lambda=\frac{u(t,\phi,\lambda,z)}{(R+z)\cos\phi},\qquad
 \dot z=w(t,\phi,\lambda,z;\text{phase})
 \label{eq:flight}
\end{equation}
を積分する。経度の距離因子には打上げ緯度ではなく\emph{現在緯度}を使う。
高度を含む半径を毎回更新する。上式はradで書いたが、実装の状態はdegreeで持ち、右辺で角速度をdegree/sへ換算して積分する。
極付近で経度が特異になる場合は停止する。現取得器は経度の継ぎ目をまたぐ箱を扱わないが、日本の任意地点の矩形取得とその範囲内の飛行には同じ式を使える。

平均風プロファイルに滞在時間を掛けて$\Delta x=\sum u_kT_k$とする方法は、風と層時間を固定する条件なら有用な比較である。
本版では風を時刻・位置で更新し、物理上昇では滞在時間も変わるため、その固定和を完全な代替式とはしない。
気候値からの平均軌道と、各日の軌道を平均したものが一般に一致するとも仮定しない。

\section{積分・イベント・途中停止の意味}\label{th:solver}
\subsection{適応刻みはSciPy、飛行の受入は本体が担う}
一般的な積分・誤差推定を重ねて独自実装する理由はないため、0.24では\code{scipy.integrate.RK45}を使う。
Dormand--Princeの5(4)次対で誤差を推定して刻みを調節し、採用区間内の状態は4次のdense outputで補間する。
各右辺$F$は、その段の時刻・位置・高度で気象を引き直す。出力点だけで風を更新する方式ではない。
本体は低levelの\code{step()}で一つずつ試行を受け取り、相、気象支持、地形、保存してよい状態を判断する。
全時間を一括の\code{solve_ivp}に渡さないのは、支持外の右辺例外でも最後の採用状態と元の停止理由を残すためである。

誤差評価の成分別スケールは
\begin{equation}
 s_i=\mathrm{atol}_i+\mathrm{rtol}\,|y_i|,
 \qquad y=(\phi_{\rm deg},\lambda_{\rm deg},z_{\rm m})
 \label{eq:error-scale}
\end{equation}
とする。\code{relative_tolerance}の既定は$10^{-9}$、\code{absolute_tolerance}は緯度・経度に$10^{-9}$ degree、高度に$10^{-3}$ mを用いる。
これは状態の単位ごとに局所誤差を制御する設定であり、全球の距離誤差や最終着地点の誤差保証ではない。
経度の原点と緯度によって角度の相対項・距離換算も変わる。高度の仮定、気象補間、イベント位置の誤差を消す設定でもない。

\code{max_step_s}は採用点間隔の上限、\code{max_steps}は採用歩数の上限である。
通常の誤差制御はSciPyへ任せる。気象の支持外や検出された途中の地形横断では、本体が失敗した試行を破棄して上限を縮め、最後の採用点から再試行する。
破裂では積分器を再初期化し、不連続な上昇・降下を一つのdense補間で結ばない。
保存時刻は適応的な採用点とイベント点であり、等時間間隔の配列ではない。固定出力時刻への再標本化は未実装である。
採用した方法、許容誤差、SciPy/NumPyの実行版は結果の\code{numerics}と保存来歴へ残す。
0.23までの自前RK4は移行比較の履歴であり、現行の選択可能backendではない。

\subsection{接地前の試行点と、保存する状態は区別する}
イベント関数$G=z-z_g(t,\phi,\lambda)$が一ステップで正から非正へ変わると、採用区間のdense output上で\code{scipy.optimize.brentq}に時間の根を求めさせる。
破裂と地形交差の候補が同じステップにあれば先のものを採り、イベント時点で相を切り替える。
\code{event_tolerance_s}は時間根の絶対許容誤差\code{xtol}であり、brentq既定の相対項も併用する。根は補間された軌道に対するもので、着地点のメートル誤差保証ではない。
指定上昇率と高度破裂の組合せでは、残り高度/速度から求めた正確な破裂時刻も使う。

イベントの高度が取得上端と一致するだけでは十分とは限らない。等温上昇などの数値的な破裂探索は、根の上側を挟む段で気象上端へ触れ、破裂直前で支持エラー停止し得る。これは旧RK4にも残った限界で、定速の解析的な破裂刻みとは異なる。取得高度には余白を持たせる。支持端を明示するadapterと有向イベントの契約を受け入れるまでは、上端外挿や破裂の推定完了で埋めない。
破裂点の上昇側記録とイベントは、到達時に先に採用する。下降側の必要量が欠ければそこまでを保って停止し、成功時は同時刻の上昇/下降2点を相で区別して保存する。

根を挟む試行のRK段が地面より下に出る場合、その段の気象照会高度だけを局所地面へ上げる。
これは終了イベントを見つけるための数学的な延長で、地下の気球状態を採用したものではない。
保存点は$z\ge z_g$を要求し、着地点は$z=z_g$へ置く。この\emph{積分段の地表固定}は、第\ref{th:weather}節の\emph{寄与列の地表延長}とは別である。
前者の回数は\code{ground_clamped_stage_queries}、後者の種類は気象品質へ残る。

\begin{center}
\begin{tikzpicture}[x=1cm,y=1cm,>=Latex,font=\small]
\path[use as bounding box](0,0) rectangle(15.5,5.4);
\node[anchor=north west,draw=accent,fill=accent!5,rounded corners=2mm,text width=6.6cm,inner sep=3mm,align=left] at(.1,5.3){\textbf{場の中：寄与列の延長}\par
照会点は混合地形より上にある。周囲の高い地形列だけが、その照会高度へ届かない。\par\medskip
\textcolor{accent}{その列の気象値だけ地表値へ延長}\par\medskip
軌道の高度は動かさない。\par
\code{terrain_stencil_surface_clamp}\par
と深さを気象品質へ残す。};
\node[anchor=north west,draw=orange!75!black,fill=orange!5,rounded corners=2mm,text width=6.6cm,inner sep=3mm,align=left] at(8.1,5.3){\textbf{積分の中：試行照会の固定}\par
接地を挟む試行状態が地面より下へ出る。根を見つける右辺の評価は続けたい。\par\medskip
\textcolor{orange!75!black}{その段の照会高度だけ地面へ固定}\par\medskip
地下の状態は保存しない。\par
\code{ground_clamped_stage_queries}\par
を計数し、着地点を地面へ置く。};
\end{tikzpicture}
\end{center}

降下中の交差だけが\code{landed}になる。上昇中に地形へ当たれば\code{TERRAIN_COLLISION}、時間・空間・高度の支持を失えば元のエラーコード、積分上限なら\code{MAX_DURATION}または\code{MAX_STEPS}を持つ\code{stopped}結果になる。
端点のほか、内部段の地表固定とdense区間中点を調べ、途中の横断を検出した区間は縮める。それでも未評価の区間だけ横切る狭い障害物や一刻み内の複数交差を完全には検出できない。細密地形と衝突の完全検出は本版の受入範囲に含まない。

\section{何を検証でき、何がまだ分からないか}\label{th:validation}
\subsection{式と数値処理は、別の期待値で確かめる}
検証は「出力が有限だった」だけで終えない。式\eqref{eq:density}は分圧による密度の和から独立に再計算し、式\eqref{eq:ascent}は出力された浮力と抗力を別に釣り合わせる。
指数降下は式\eqref{eq:exponential}の解析時間、一定密度・東向き風は$\int u/(R+z)dt$の対数積分、一定$u/v$の斜行は$\lambda=(u/v)\ln(\sec\phi+\tan\phi)+C$と比較する。
これらは実装中の同じ関数をもう一度呼んで一致を得る検査とは異なる。

\begin{longtable}{@{}p{.25\linewidth}p{.36\linewidth}p{.33\linewidth}@{}}
\toprule 問い & 独立な期待・反例 & 対応する実装/試験 \\\midrule\endhead
補間順序が正しいか & 時刻で気圧面高度が変わる場、アフィン場、端点と欠測 & \path{WeatherField.sample} / \path{tests/test_flight_weather.py} \\
物理量の定義が正しいか & 分圧恒等式、浮力抗力釣合い、非正自由揚力 & \path{moist_air_density}, \path{vertical_state} / \path{tests/test_trajectory.py} \\
全飛行が解けるか & 解析飛行時間、球面の積分、斜面接地、直径破裂 & \path{simulate} / 同試験 \\
支持不足を成功にしないか & 途中の時刻・領域境界、地表下放球、最大時間/歩数 & \path{simulate} / 同試験 \\
入力を追跡して再実行できるか & rawハッシュ改変拒否、同版の再復号、設定/入力/コードのハッシュ & \path{replay_gfs}, CLI保存 / 気象・CLI試験 \\\bottomrule
\end{longtable}

最大刻みを半分にした比較に加え、成分別の許容誤差を厳しくしたときの着地時刻・位置・最大高度の差を確認する。適応刻みでは保存点数も変わるため、同じ行番号を同じ時刻として比較しない。単一条件の一致から収束次数や現実の精度は証明できないが、数値設定への感受性を測れる。
北海道の簡易例と和歌山の等温例は、本版の実GFS保存場を入力した受入例である。最終実行値・条件・刻み比較・再現ハッシュは実施記録と結果JSONを正本とし、本文の一般的な精度として固定しない。

\subsection{提供された風解析から何を受け入れたか}
ユーザー提供の『JRA-3Qを用いた高層風の統計的傾向の解析』p12、p88--89、p97--99は、風の成分扱い、時間/場所固定の効果分離、地図因子、気候値と予報誤差の区別を具体的に説明している。
その問いと比較の組立てを本版へ生かした。一方、提供された北海道DBは半月統計であり、そのまま瞬時の気象場にはしない。
独立レビューでは全519840要約行の代数関係を調べ、実体の上端3 hPa、全期間14612時刻と、説明中の不整合を記録した。
有効標本数には記載式にない上限処理が見られた。原JRA時系列と和歌山の原統計は未提供であり、報告の誤差kmや改善率を実証済みの性能として受け入れてはいない。

同報告p103--109の層間相関の論点は、将来の不確実性評価に効く。各高度の周辺分布を独立に乱数化するだけでは、風速・風向・高度・時間・空間の依存を失い得る。
単一飛行核は決定論的である。0.48で加える入力の仮分布も、気象場を無条件に乱数化するものではない。標本内の経験的な被覆を、校正した着地確率や安全な落下域と区別する。

\subsection{次のモデルを足すときに守る境界}
初版の等温仮定は、熱モデルを永久に排除する判断ではない。熱を加えるなら$T_g$等を状態へ加え、エネルギー式と放射/伝熱入力、適用範囲、単純極限への戻りを同じ変更で記述する。
慣性を加えるなら$w$を状態にし、準定常速度を代入する式と加速度の式を混ぜない。
Re依存抗力なら$\mu,T,\rho,d,w$の依存を明記し、熱と抗力の変更が同時に有効になった場合も効果を分離できる比較を残す。
この拡張は追加実験を本体着手の必須条件にするものではなく、得られた根拠の範囲で比較候補を実装し、精度改善を主張する段階で別に評価する。

\clearpage
\section{同じ標本で条件を比べ、干渉群へ戻る}\label{th:ensemble}
\subsection{最初に固定するのは入力の仮説である}
例えば「ガスの充填量をこの幅で変えたとき、30分延期すると干渉する標本がどう変わるか」を考える。
ここでの上下限は利用者が理由とともに指定する感度仮説であり、製品のばらつきや予報誤差を観測から推定した値ではない。
初回は元量を一つ選び、
\begin{equation}
 X_j=a+(b-a)U_j,\qquad U_j\sim\mathcal U[0,1),\quad j=1,\ldots,N
 \label{eq:sensitivity}
\end{equation}
とする。NumPyのPCG64と一様乱数を利用し、方式・実行環境・seedに加え、実際の$X_j$、単位、重み$w_j=1$、標本IDを保存する。
再表示や再試行で再抽選しない。乱数ライブラリの版を越える再生成保証をseedだけに求めない。

\begin{longtable}{@{}p{.28\linewidth}p{.29\linewidth}p{.37\linewidth}@{}}
\toprule 選べる元量 & 使うモデル & 保つ関係 \\\midrule\endhead
ガス質量 [kg] & 内外等温の浮力・抗力 & 同じ値をガス体積と上昇質量へ使う。別に指定した破裂後有効質量へ機械的に足さない \\
破裂高度 [m ASL] & 高度による破裂 & 放球高度を越えるという既存の条件を守る \\
定格下降速度 [m/s] & 密度補正した定格下降 & 基準密度等の固定値を保持し、選んだ速度だけを変える \\\bottomrule
\end{longtable}
選んだ量を使わないモデルは、効果がないまま感度試験と呼ばず、計画で明示拒否する。
個々の実現値が物理入力として不正なら、その標本と理由を台帳へ残す。別の値に引き直す・境界へ丸める操作は、指定分布を変えるため行わない。
荷重、自由揚力、製品の破裂径・膜厚は、関係する状態を一緒に解決する設計が必要であり、この三つの元量へ名前だけを付け替えて追加しない。

\begin{center}
\begin{tikzpicture}[node distance=5mm, box/.style={draw=accent,rounded corners,fill=accent!6,text width=.84\linewidth,align=left,inner sep=3mm},arr/.style={-{Stealth},thick,accent}]
\node[box] (draw) {\textbf{一度だけ抽選}：元量$X_j$、標本ID、単位、仮定理由を固定};
\node[box,below=of draw] (cases) {\textbf{同じ物理値を各候補へ}：親、延期、対応するモデルへ完全な入力を解決};
\node[box,below=of cases] (runs) {\textbf{全予定試行を保存}：着地、物理停止、未実行、失敗を区別};
\node[box,below=of runs] (analysis) {\textbf{同じ結果から読み直す}：経験範囲、禁止域の群、入力と履歴、対応する変化};
\draw[arr] (draw)--(cases);\draw[arr] (cases)--(runs);\draw[arr] (runs)--(analysis);
\end{tikzpicture}
\end{center}
延期は放球時刻だけを変え、同じ$X_j$と同じ固定気象資産で比較する。元量を両モデルが使う場合はモデル比較にも対応付けられる。
分布の中心・幅を変更した計画は別の標本集合である。同じ番号や同じseedだけで、同じ物理標本だと判定しない。
気象の実現値IDは機体の標本IDと分ける。過去数年を代表する時期・時間帯の抽出規則や、場の相関を保った不確実性は後続の別規則として加える。次節の明示原日時入力とは区別する。


\subsection{原日時の場を明示して比べる場合}\label{th:historical-windows}
0.55の過去窓経路では機体の乱数を新たに抽選しない。共通機体条件$\theta_c$と、
利用者が選んだ原放球日時$t_i$・その日時を支える場$W_i(\boldsymbol{x},z,t)$を固定し、
\begin{equation}
 Y_{c,i}=\mathcal F(\theta_c,t_i,W_i),\qquad i=1,\ldots,N
\end{equation}
を同じ飛行核で求める。$\mathcal F$の時間・鉛直内挿は前章の場の定義のままである。
各$t_i$は秒まで保持し、同じ原UTCを別の名称で重複計上しない。窓IDと、場の内容hash・原有効UTC、
選定理由を残す。原日時を所有する台帳と、予報の親子延期指定を同時に適用しない。

場が未取得、時間支持が足りない、放球点が領域外などの行も予定$N$に残す。
実飛行は支持された行のみで行う。経験範囲は着地した標本（着地数$L$）、相別履歴は着地と途中停止を含む有効な原記録を対象とし、各経過時刻の支持数を保つ。
重みは各行1である。例えば冬と夏の二日だけを選んだ結果なら、示されるのはその二日に条件付けた比較であり、
冬と夏が同じ頻度で生じるという母集団モデルではない。$L=2$では経験域は線へ退化し得る。
未取得の多い季節を、着地例が少ないから安全と読むこともできない。

半月集計から渡す地域・季節・資料指紋は\emph{何を調べたかったか}の記録である。
集計済み風から高度・空間・時間の同時相関を再生成した場ではない。
将来の季節計画では、対象年・月内日・時間帯の母集団、抽出法、取得網羅率、
機体の不確実性との組合せを別途定める必要がある。この有限接続はその議論を実計算へ結ぶ入口であり、
二つの診断日から多年の将来確率を受け入れるものではない。

\clearpage
\subsection{予定数、着地数、干渉数は違う分母である}
一つの候補の予定数を$N$、着地数を$L$、禁止領域の和集合に着地した数を$H$とする。
概要で読む$L/N$はどれだけ着地まで求まったか、$H/L$は求まった着地のうち何本が干渉したかを表す。$L=0$なら$H/L$は未定義であり、0\%とはしない。
物理停止と、計算失敗・取消・中断・未実行を分ける。履歴が空の初期不成立も全予定台帳から消さず、再試行attemptを新しい機体標本として数えない。

禁止域との判定は\emph{すべての着地点}に行う。95\%経験域の外の一点も対象であり、経験楕円と領域が見た目で交わることを干渉数の代用にしない。
現画面のPolygon/MultiPolygon判定は、外周・穴の境界を含み、穴の内部を除外する。領域間の重複は和集合として一度だけ数え、施設別の名前と群も保持する。
この判定は画面側の固定された領域・実装による分類である。分析APIは渡された標本IDの集計を保存し、サーバー側が独立に領域判定を認証したとは扱わない。

\subsection{50・90・95\%は標本の経験範囲である}
着地点をWGS84の局所方位距離平面へ写し、東・北の二次元ベクトル$\boldsymbol x_i$ [m]とする。
投影の中心は着地点の緯度・経度の算術平均で、表示する平均位置は平面上の平均を逆変換する。
現実装は極域・日付境界を跨ぐ描画、および中心から1000 kmを越える描画を未対応とする。この表示制限によって着地数や干渉判定の標本を減らさない。

\begin{equation}
 \overline{\boldsymbol x}=\frac1L\sum_i\boldsymbol x_i,\qquad
 C=\frac1{L-1}\sum_i(\boldsymbol x_i-\overline{\boldsymbol x})(\boldsymbol x_i-\overline{\boldsymbol x})^\mathsf T
\end{equation}
$L=1$は一点とする。$C$をNumPyの対称行列固有分解で扱い、支持された固有方向だけで二乗半径
\begin{equation}
 d_i^2=(\boldsymbol x_i-\overline{\boldsymbol x})^\mathsf T C^+
       (\boldsymbol x_i-\overline{\boldsymbol x})
\end{equation}
を求める。$C^+$は支持方向での逆行列である。
目標割合$p\in\{0.50,0.90,0.95\}$に対し、昇順の第$k=\lceil pL\rceil$番目を$s_p$とし、$d_i^2\le s_p$の標本を含む領域を描く。
これは正規分布の確率楕円や平均の信頼領域ではない。計算された半径が同じ標本はすべて含み、実際の含有数$K_p$と$K_p/L$を指定割合とは別に保持する。
小さい$L$や同順位では50・90・95\%が同じ輪郭になる。それを広げて異なる面に見せない。

二方向が支持されれば楕円、一方向なら線分、ゼロ方向なら点で示す。
数値退化とするのは捨てる方向の最大残差が$10^{-6}$ m以下の場合だけであり、細い分布へ人工的な横幅を足さない。負の固有値等で解像できない場合は算出不可を示す。
楕円半径の順位と、地図へ描く有限個の頂点による近似は別である。境界上の観測方向を描画頂点にも含め、投影前後の丸めや描画近似を科学的な精度保証へ使わない。

\textbf{読み違えやすい場面：}曲がった軌道族や二峰の着地では、楕円が標本の少ない中間領域も塗る。
「塗られたどこにも同じ確率がある」「空白より危険である」とは読めない。必要なら点と群を掘り下げる。
集計APIは選ばれた着地点の凸包もSciPyで求め、点・線への退化を保持する。ただし概要・3D・KMLの干渉表示は個別着地点とし、件数から該当ID群の詳細へ進む。凸包を干渉面や100\%の将来確率領域として表示しない。

\clearpage
\subsection{履歴は共通経過時刻と相で揃える}
各飛行の不等間隔な原記録を保持したまま、描画用の共通時刻$t$で同じ相の隣接記録を線形内挿する。
基本の時刻列は最大241点であり、各相の始端・終端を追加する。この表示用時刻列を数値積分の保存刻みと混同しない。
時刻$t$・相$r$で内挿できる標本集合を$A_r(t)$とし、量$y$について
\begin{equation}
 \mu_r(t)=\frac1{|A_r(t)|}\sum_{j\in A_r(t)}y_{j,r}(t)
\end{equation}
と10・90分位の帯を描く。NumPyの線形標本分位を使い、各量の有効数も示す。集合が空なら値は欠測とし、0で埋めない。
高度、東北変位、累積横移動距離、水平・鉛直速度がこの対象である。東北変位は飛行核と同じ半径の球面で放球点からの距離・方位を成分化し、累積距離は隣接保存点の大円弧長を加算する。

破裂時刻の上昇側の値と下降側の値は各相の線端として別々に含める。この瞬間の相別有効数を合算しない。
着地・停止時刻の最終値は端点として描き、その\emph{後}へ値を延長しない。瞬間の風分布では着地済み標本を除外し、履歴図の端点規約と区別する。
放球直後の破裂で同じ相・時刻の記録が二つあるとき、物理値が同じなら分析用に一つへまとめ、矛盾する場合は拒否する。原記録やイベントは消さない。

平均位置を結んだ線の構成員は時刻・相で変わる。その線自体が一本の実在する気球の軌道ではない。
干渉群と全着地の入力分布・履歴が違っても、それだけで因果的な主因を識別したとはしない。低ガス量で浮上できなかった標本等を調べるため、予定$N$の実現値と状態へも戻れるようにする。
変更案は選択した干渉群だけでなく\emph{全元標本}を再評価する。同じ標本の干渉退出・残留・新規流入を区別し、片側未着地・未計算を差ゼロで埋めない。

\subsection{有限な検証と、後続へ残す問い}
同点、直線、薄い面、二峰、ゼロ着地、空の停止履歴、異なる破裂時刻、即時破裂、取消・再試行を反例とする。
数値試験は母数と幾何・時刻の意味を守るための検査で、製品分布、予報の校正、実機の着地精度を実証するものではない。
保存は固定計画、全予定行、実行attempt、不変の結果snapshotを分ける。図の選択をsnapshot・候補・領域版・標本IDへ結び、再試行で前の図を変えない。
今後、重み付きの過去年、複数元量、気象場の実現値を加えるときは、選択母集団・相関・重みと計算費用を先に定義し、現在の等重み式へ無言で混ぜない。


\clearpage
\section{過去の風を、時期・高さ・地域で比べる}\label{th:climate}
\subsection{答える問いと、資料ごとの母集団}
「どの時期の、どの高さで風が強いか」「向きがそろうか」「地域内の違いは大きいか」を、年間の高度別図・地域図・月の中央値と分位幅で読む。平均プロファイルは補助図である。代表地点の風配は、その地点の経験頻度を示す。
これは長期計画の候補を絞る入口である。平均風が弱い時期を見つけても、個々の飛行の移動距離や落下分散が小さいとはまだ結論できない。

資料は二系統ある。提供されたJRA-3Q統計DBは2016--2025年の保存平均と件数を持つ。独立した原標本束は、指定年のUTCごとの東西/南北風$u,v$を持つ。どちらもUTC 00・06・12・18時から選び、各月の1--15日と16日--月末の\textbf{24半月}を基礎とする。原標本からの年間図だけには、各半月を二分した48区分も別に用意する。
\begin{center}\small
\begin{tabular}{@{}p{43mm}p{45mm}p{67mm}@{}}
\toprule 資料 & 現候補の固定支持 & 読める分布と限界 \\\midrule
提供JRA統計DB & 2016--25年・570格子。全時刻38面、時刻subset20面。 & 元半月/地点統計を結合。原標本なしで月/地域分位を再構成しない。\\
JRA原標本の例 & 2024年・38面・570 Gaussian格子 & 同じ元UTCのu/vから選択領域の平均、経験分位、代表点風配。\\
ERA5原標本の例 & 2024年・37面・1188等緯度経度格子 & native座標/面を保持。JRAへ補間せず、年も旧DBへ延長しない。\\\bottomrule
\end{tabular}\end{center}
上記は実装対象の支持であり、取得・実画面の完了実績はS36で確認する。年は資料ごとに固定し、画面で年の部分選択やJST暦への再区切りは行わない。2024年の図を多年の代表性や来年の予測分布と呼ばない。

\begin{center}
\begin{tikzpicture}[font=\small,>=Latex]
\node[draw=accent,fill=accent!7,rounded corners,text width=43mm,minimum height=16mm,align=center] (a) {半月・月・季節 $g$\,$\times$\,面 $k$\\格子内で時間を結合};
\node[draw=accent,fill=accent!7,rounded corners,text width=43mm,minimum height=16mm,align=center,right=8mm of a] (b) {選んだ格子集合 $I$\\提供格子重みで集約};
\node[draw=ink,fill=ink!5,rounded corners,text width=43mm,minimum height=16mm,align=center,right=8mm of b] (c) {年間・地域・季節の比較\\平均風速、向き、定常度};
\draw[->,thick,accent] (a)--(b);\draw[->,thick,accent] (b)--(c);
\end{tikzpicture}
\end{center}

一つの格子・半月・面について、時刻数を$n_b$、保存された平均東西風・南北風・スカラー風速を$\bar u_{bki},\bar v_{bki},\bar s_{bki}$と書く。
意味としては、同じ時刻集合$T_b$に対する
\begin{equation}
 (\bar u_{bki},\bar v_{bki},\bar s_{bki})
 =\frac{1}{n_b}\sum_{t\in T_b}
 \bigl(u_{tki},v_{tki},\sqrt{u_{tki}^2+v_{tki}^2}\bigr)
 \label{eq:climate-time}
\end{equation}
である。提供DBの経路は\emph{保存済みの三つの平均と件数}を読む。原標本経路は同じUTC集合のu/vから速さを計算し、この三つの平均を作る。二系統を同じ観測期間と取り違えない。

\begin{center}\small
\begin{tabular}{@{}lrrp{65mm}@{}}
\toprule UTC半月 & 10年間の日数 & $n_b$ & 件数から言えること \\\midrule
1月1--15日 & 150 & 600 & 同じ半月・面で、全対象格子の件数一致を要求。\\
1月16--31日 & 160 & 640 & 月長によって時刻数が変わる。\\
2月16--月末 & 133 & 532 & 2016・2020・2024年の閏日を含む。\\\bottomrule
\end{tabular}
\end{center}
この表は旧DBの例で、10年間・24区分の和は14,612時刻。2024年の原標本全4時刻は366日で1,464時刻となる。\textbf{時間の件数$n_b$と空間の格子数$|I|$は別に表示する。}
$n_b|I|$を独立な天候の標本数にしない。件数一致だけでは同一時刻集合や独立性を証明できないため、旧集約DBの元時刻の対応と有効独立標本数は未確認である。新原標本束では暦上のUTC列を照合するが、有効独立標本数を算出したことにはならない。

\clearpage
\subsection{半月を月・季節へ結び、対応時刻を選ぶ}\label{th:climate-periods}
期間$g$の各格子で、互いに重ならない元区分$j$の件数を重みにする。全時刻なら$j$は半月、時刻選択なら「半月と選択UTC時刻」の組である。$x$を$u,v,s$のいずれかとして
\begin{equation}
 N_g=\sum_{j\in g} n_j,\qquad
 \bar x_{gki}=\frac{\sum_{j\in g} n_j\bar x_{jki}}{N_g}.
 \label{eq:climate-pool-time}
\end{equation}
日数の異なる半月を等重みで平均しない。全格子の$N_g$と暦上の期待件数を照合した後に、後述の「地域平均は、成分と速さを別々に集約する」へ進む。
DJF/MAM/JJA/SONは12--2月/3--5月/6--8月/9--11月のpoolである。\textbf{DJFは各資料の指定年内の12月・1月・2月の集合}であり、年をまたぐ完全な冬が連続しているとの意味ではない。

\begin{center}\small
\begin{tabular}{@{}lrrrr@{}}
\toprule 同じ解析時刻 & \multicolumn{4}{c}{提供された4時刻の対応} \\\midrule
UTC & 00 & 06 & 12 & 18 \\
JST & 09 & 15 & 21 & 翌日03 \\\bottomrule
\end{tabular}
\end{center}
JSTは時刻表示の換算であり、期間の所属日はUTCのまま。$N_g$は期間内日数と選択時刻数の積となる。提供DBの時刻別平均は\textbf{1000--300 hPaの20面}、全4時刻は38面。原標本束はその資料の全支持面で時刻を選べる。提供時刻を含まない選択、未対応面、欠けた支持をゼロや別母集団へ置き換えない。

\paragraph{年間図だけに使う月内4区分}
原標本では各月$m$の日を\textbf{1--7日、8--15日、16--22日、23日--月末}に分ける。月末$D_m$は資料年の実暦に従い、閏日を含める。区分$q$に属する選択UTC時刻の集合$T_{mq}$に対し、式\eqref{eq:climate-time}の和を元u/vから直接計算する。件数は各区分の日数と採用UTC時刻数の積で、地域集約には従来と同じ格子重みを使う。ISO週やJST日付ではなく、1--15日/16日--月末という旧半月境界の内部を分ける定義である。
年間の48区分は季節内の変化を見るための表示で、月中央値2本と月全体p10--p90の母集団は変えない。地域図・風配・飛行関心には従来の半月/月/季節を使う。提供10年DBの半月平均からこの4区分を逆算せず、同じDBの再提供も元時刻の支持を増やさない。2024年の原標本を細かく表示しても、10年の代表性や独立な天候の標本数は増えない。

\clearpage
\subsection{月の中央値と幅は、元の格子・時刻をまとめて求める}\label{th:climate-quantiles}
問いは「前半と後半で典型的な速さがどう変わり、月内にはどれほどの幅があるか」である。各面$k$で、選択地域の元格子$I$と月$m$の選択UTC時刻$T_m$を組にする。風速$s_{tki}=\sqrt{u_{tki}^2+v_{tki}^2}$、正規化した格子面積重み$w_i$を用いて、
\begin{equation}
 F_{mk}(x)=\frac{1}{|T_m|}\sum_{t\in T_m}\sum_{i\in I}w_i\,\mathbf{1}[s_{tki}\le x],\qquad
 Q_{mk}(q)=\inf\{x:F_{mk}(x)\ge q\}.
 \label{eq:climate-quantiles}
\end{equation}
各UTC時刻は等重み、各格子は面積重み。格子を先に平均した時系列$\sum_iw_i s_{tki}$の分布ではない。前半$T_{m,1}$・後半$T_{m,2}$の中央値はそれぞれ同じ式で別々に求め、月の帯は全$T_m$の$Q(0.1)$と$Q(0.9)$を使う。二つの半月分位や地点分位を平均して月/地域分位にしない。

\begin{center}
\begin{tikzpicture}[font=\small,>=Latex]
\node[draw=accent,fill=accent!7,rounded corners,text width=40mm,align=center] (a) {選択した元格子$\times$UTC\\面ごとの風速と重み};
\node[draw=accent,fill=accent!7,rounded corners,text width=48mm,align=center,right=7mm of a] (b) {前半：$Q_{m,1}(0.5)$\\後半：$Q_{m,2}(0.5)$\\月全体：$Q_m(0.1),Q_m(0.9)$};
\node[draw=ink,fill=ink!5,rounded corners,text width=43mm,align=center,right=7mm of b] (c) {12か月を共通の軸で比較\\青緑/橙：中央値2線\\灰帯：月10--90\%};
\draw[->,thick,accent](a)--(b);\draw[->,thick,accent](b)--(c);
\end{tikzpicture}\end{center}

実装はNumPyの\code{quantile(method='inverted_cdf', weights=...)}で式\eqref{eq:climate-quantiles}の加重経験CDF逆関数を評価する。人工模型の等重み・線形補間分位と有限標本の定義が異なる。例えば風速2/10 m/s、重み1:3の2値なら、$F(2)=0.25$で$Q(0.1)=2$、$Q(0.5)=Q(0.9)=10$となる。これは定義の説明例で、実データの結果ではない。

欠測を除外して別母集団にしない。取得済み月は指定年の全UTC列、全支持面・全格子を照合し、未取得月は空欄として残す。独立raw資料は指定年の全12か月を必要とする。旧DBに同支持raw月束を加える経路だけは、未取得月の分位を欠測として示せる。図の幅は経験的な風速の範囲であり、中央値の信頼区間、将来の予測区間、飛行の安全確率ではない。

\subsection{代表地点の風配は、件数を合算する}\label{th:climate-rose}
旧DBの提供風配は固定元格子285（北緯43.2674度、東経143.6250度）の全4時刻である。地域を狭めても地点は移動しない。選択地域の内外を明示し、地域全体の風向・風速分布とは呼ばない。
半月$b$・面$k$・方向階級$d$・速度階級$c$の件数$C_{bkdc}$から
\begin{equation}
 C_{gkdc}=\sum_{b\in g}C_{bkdc},\qquad
 f_{gkdc}=\frac{C_{gkdc}}{N_g},\qquad
 \sum_{d,c} C_{gkdc}=N_g
 \label{eq:climate-rose}
\end{equation}
を作る。分母はその一点の気象解析時刻数であり、格子数を掛けない。頻度や月の分位を等重みで平均して作る図ではない。
来向FROMは北0度・時計回りの16方位。幅22.5度で北の階級は0度をまたぐ。速度は$[0,5),[5,10),[10,20),[20,30),[30,50),[50,\infty)$ m/s。0--5 m/sも通常階級で、静穏の独立区分はない。
\textbf{旧DBには時刻を絞った風配の元件数がない}ため、その選択を全時刻へ黙って置き換えない。

原標本束は既定では資料中心に最も近い元格子を使い、取得済みの別の元格子も選択できる。その一点・選択UTCのu/vから同じ階級を数え、地域を動かしても点は移動させない。地域の分位と一点の風配を混同しない。数値上$s=0$の静穏は方向を持たないため、方向階級に入れず$C_{\rm calm}$へ置く。分母は静穏を含む全$N_g$のままとし、
\begin{equation}
 C_{\rm calm}+\sum_{d,c}C_{gkdc}=N_g,\qquad f_{gkdc}=C_{gkdc}/N_g
\end{equation}
を検査する。静穏数が正なら$n/N$を図に添える。動いている標本だけで100\%に再正規化せず、偽の北風も作らない。旧DBの「静穏区分なし」という定義を後から書き換えない。

\clearpage
\subsection{地域平均は、成分と速さを別々に集約する}
資料の正の格子面積重み$a_i$を、選択した格子中心の集合$I$内で正規化する。JRAはGaussian格子の重み、ERA5の等間隔緯度経度格子は緯度帯の球面面積（この等間隔格子では$\cos\phi_i$に比例）を使う。
ここでのGaussianは格子の重みであり、風が正規分布に従うという仮定ではない。矩形の端にかかるセルの面積割合を積分する方式でもない。
\begin{equation}
 w_i=\frac{a_i}{\sum_{j\in I}a_j},\qquad
 U_{gk}=\sum_{i\in I}w_i\bar u_{gki},\quad
 V_{gk}=\sum_{i\in I}w_i\bar v_{gki},\quad
 S_{gk}=\sum_{i\in I}w_i\bar s_{gki}.
 \label{eq:climate-space}
\end{equation}
$S$は風速の平均、$\sqrt{U^2+V^2}$は平均風ベクトルの長さである。向きの異なる風が相殺されるため、両者は一般に一致しない。
年間図と平均プロファイルは地域集約の$S$を使い、年間図の来向は$(U,V)$から得る。地域図では各格子の$\bar s_i$を色、$(\bar u_i,\bar v_i)$の移流先TOを等長矢印にする。矢印の長さを風速と読まない。

\textbf{地域で集約した定常度}を
\begin{equation}
 R_{\mathrm{pool},gk}=\frac{\sqrt{U_{gk}^2+V_{gk}^2}}{S_{gk}}
 \quad(S_{gk}>0),\qquad
 R_{\mathrm{local},gki}=\frac{\sqrt{\bar u_{gki}^2+\bar v_{gki}^2}}{\bar s_{gki}}
 \quad(\bar s_{gki}>0)
 \label{eq:climate-constancy}
\end{equation}
とする。地域図は各格子の$R_{\mathrm{local}}$、年間の地域集約図は$R_{\mathrm{pool}}$を示す。
前者はその格子での時間的な方向の相殺、後者はそれに加えて\emph{格子間の方向の相殺}も含む。
地域の低い$R_{\mathrm{pool}}$だけを見て、各地点の風向が時間的に不安定だと断定しない。

\begin{center}\small
\begin{tabular}{@{}p{49mm}rrrr@{}}
\toprule 同じ速さで逆向きの2格子 & $U$ & $S$ & $R_{\rm pool}$ & 局所$R$の加重平均 \\\midrule
$\bar u_1=+10,\ \bar u_2=-10$、重み1:1 & 0 & 10 & 0 & 1 \\
同じ2格子、重み1:3 & $-5$ & 10 & 0.5 & 1 \\\bottomrule
\end{tabular}
\end{center}
両格子とも$\bar v=0$、風速はm/s。\textbf{局所$R$を平均しても、地域で集約した$R$にはならない。}
$R$は方向のそろい方を速さとともに要約する量で、風速の変動係数や予報誤差ではない。

\begin{center}\small
\begin{tabular}{@{}lrrrrr@{}}
\toprule 全570格子・300 hPa & $U$ [m/s] & $V$ [m/s] & $S$ [m/s] & $R_{\rm pool}$ & $n_b$ \\\midrule
1月前半 & 31.594 & $-0.452$ & 35.017 & 0.902 & 600 \\
7月前半 & 15.455 & $-0.805$ & 20.839 & 0.743 & 600 \\\bottomrule
\end{tabular}
\end{center}
従来の半月接続で、提供平均を別経路の加重和と照合した例を継承する。ここからは当該面の季節差を読めるが、風が弱い半月の着地精度を保証しない。
地域内の風速幅$\max_i\bar s_{gki}-\min_i\bar s_{gki}$は、時間結合後に再計算する。半月の幅を平均しない。これは\emph{格子平均間の幅}であって、個々の時刻の風速分布や分位幅ではない。

\clearpage
\subsection{向きが未定義のときと、高さを読むとき}
FROMは北0度・時計回りの「吹いてくる方位」である。平均成分から最後に
\begin{equation}
 \theta_{\rm FROM}=\left[\operatorname{atan2}(U,V)\frac{180}{\pi}+180\right]\bmod360
 \quad\bigl(\sqrt{U^2+V^2}>0\bigr)
 \label{eq:climate-from}
\end{equation}
を作る。局所表示では$U,V$を$\bar u_i,\bar v_i$へ置き換える。方向角そのものを平均しない。
\begin{center}\small
\begin{tabular}{@{}p{40mm}p{39mm}p{77mm}@{}}
\toprule 平均の状態 & 表す値 & 読み方 \\\midrule
$S=0$、$U=V=0$ & $R$とFROMは未定義 & 静穏。未定義を北風0度として塗らない。\\
$S>0$、$U=V=0$ & $R=0$、FROM未定義 & 動いていても向きが完全に相殺される。\\
平均ベクトルが非ゼロ & $R$とFROMを表示 & 小さい$R$の向きは、全ての風を代表しない。\\\bottomrule
\end{tabular}
\end{center}
理論上は三角不等式から$0\leq R\leq1$となる。実装は保存FLOATの丸めに限って$\sqrt{U^2+V^2}-S\leq10^{-6}\max(1\,\mathrm{m/s},S)$を許し、上端を1に丸める。
この範囲を超える矛盾、非有限値、欠測・件数不一致を黙って除外して平均しない。集計のゼロ判定は実装上のゼロであり、弱風を一律に静穏へ置く閾値は導入していない。

データの面は\textbf{気圧面}で識別する。年間図の主軸は固定ISA標準大気の概算高度[km]を線形に取り、気圧の副目盛[hPa]を添える。
同じ気圧面の実高度は時刻・場所で変わる。提供資料の幾何高度の全域・全期間要約を、選択地域のその時刻の実高度と置き換えない。
新原標本で取得しているのはu/vだけで、幾何高度は未取得である。ISA参考高度を実高度として補わない。1000 hPaは地表ではなく、1--3 hPaまでの統計があっても地上から連続した飛行場が得られたことにはならない。

\subsection{この接続で採る定義と、次に必要な科学的判断}
\begin{center}\small
\begin{tabular}{@{}p{27mm}p{131mm}@{}}
\toprule 区別 & 本節での扱い \\\midrule
採る定義 & UTC期間のn重み結合、等しい時間支持の検査、資料別面積重み、加重経験CDFの逆関数、集約後のR/来向、元階級と静穏のcount/分母。未定義値を区別する。\\
今回の支持 & 旧DBと原標本束を別sourceにする。年・格子・面・UTC支持は冒頭表のとおり。半月/月/季節の平均と風配、12か月の中央値/幅。原標本の年間図だけ48区分を追加。年選択・任意地点補間は含まない。\\
後続の科学検討 & 年別や未提供時刻の支持、気象の空間・鉛直・時間相関、必要な原場の取得と誤差評価。供給資料の能力と今回表示の範囲は分ける。\\\bottomrule
\end{tabular}
\end{center}
半月ごとの周辺統計や分位だけから、同時刻の共同気象場は復元できない。気圧面ごとに独立な乱数を振ると、層間相関や場の連続性を失う。
原標本に由来する今回の月分位も、各面の周辺分布である。元分位の平均を月・季節や地域の分位にすることはできない。\textbf{本節の図を、そのまま飛行用気象場や落下確率へ転用しない。}
原再解析への値の照合、実飛行に対する精度、本番図の実UI性能は今回の有限な式・DB照合とは別の受入事項として残る。

\clearpage
\section{参照と更新の入口}\label{th:references}
この文書の式は本版が実装する仮定を導出したもの、設定値はコードと利用者入力、実測は実施記録として区別する。
気象製品の仕様・取得制約は\path{docs/WEATHER_DATA_GUIDE.md}、場の契約とモデル候補は\path{docs/ENVIRONMENT_CONTRACT.md}、採否理由は\path{docs/DECISIONS.md}のD-139、D-148、D-151、D-154、構造修復のD-162、数値移行のD-163を参照する。
手順と検査実績は\path{BOOTSTRAP_RUNBOOK.html}、出典受入範囲は\path{REFERENCE_LOG.md}が担当する。
提供風報告の原本ハッシュ、ページ別の受入判断、独立計算との不一致は\path{references/WIND_REPORT_REVIEW.md}を参照する。
第\ref{th:climate}節は\path{backend/climate/}の\path{periods.py}（UTC区分）、\path{aggregation.py}（時間/地域の結合）、\path{statistics.py}（R/来向）、\path{extensions.py}（旧DBの対応時刻/経験風配）、\path{samples.py}（加重月分位）、\path{archive.py}（原UTC平均/風配とISA参考高度）に対応する。\path{source.py}の支持・原本検査とDB/API/固定保存は\path{docs/IMPLEMENTATION_NOTES.md}へ分ける。今回の追補だけで旧節の科学的採用範囲は広げない。
0.48の入力感度と固定結果の採否はD-186/S36の継続、具体的な接続はINTEGRATION第7.3節を読む。乱数・分位・固有分解・凸包・測地計算は既存ライブラリを使い、今回独自に決めたのは母集団、入力の意味、相と選択の規約である。
\begin{itemize}
\item NumPy Generator、eigh、quantile：\url{https://numpy.org/doc/stable/reference/random/generator.html}、\url{https://numpy.org/doc/stable/reference/generated/numpy.linalg.eigh.html}、\url{https://numpy.org/doc/stable/reference/generated/numpy.quantile.html}。
\item SciPy ConvexHull：\url{https://docs.scipy.org/doc/scipy/reference/generated/scipy.spatial.ConvexHull.html}。
\item GeographicLib 2.1のDirect/Inverse：\url{https://geographiclib.sourceforge.io/html/python/code.html}。
\end{itemize}
以上は2026年9月29日に一次文書を確認した。実行に固定した版はNumPy 2.4.4、SciPy 1.17.1、GeographicLib 2.1であり、stable URLの表示版と混同しない。

第\ref{th:display-height}小節の表示基準は、2026年10月2日に\href{https://docs.ogc.org/is/12-007r2/12-007r2.html}{OGC KML仕様}のabsolute高度と、\href{https://cesium.com/learn/cesiumjs/ref-doc/Cartesian3.html#fromDegrees}{Cesiumの楕円体高}を照合した。
EGM96の原配布元は\href{https://earth-info.nga.mil/php/download.php?file=egm-96interpolation}{NGAの15分補間格子}、同梱DACは\href{https://github.com/CesiumGS/cesium-native/blob/1bca8021711294cc83a5a032599ac72d9b3166c9/data/WW15MGH.DAC}{Cesium Nativeの固定commit配布物}である。
全点の原GRDとの量子化差、配布物のSHA-256、Public Domainの根拠、参照試験と独立実装の範囲は\path{frontend/scene3d/data/README.md}に固定する。
格子点・補間値・符号の照合は表示基準の確認であり、写真地物や実飛行の高度精度を受け入れた証拠ではない。


式\eqref{eq:density}の混合気体の前提と比湿の定義は、ECMWFの\emph{Atmospheric thermodynamics}（2015年、本文p1の式3・8、p4の第6節、p5の式23--25）と照合した。
\url{https://www.ecmwf.int/sites/default/files/elibrary/2015/16954-atmospheric-thermodynamics.pdf}。
本実装の$R_d=287.05$、$R_v=461.5$ J/(kg K)は同資料の丸め値287.06、461.52と少し異なり、定数まで原資料の複写だとはしない。

抗力の面積・係数の定義はNASAの\emph{Drag Equation}、等温・内外同圧・球形を置く気球の基本関係はNASAの\emph{Designing a High Altitude Balloon}を参照した。
\url{https://www1.grc.nasa.gov/beginners-guide-to-aeronautics/drag-equation/}、
\url{https://www.grc.nasa.gov/www/k-12/Numbers/Math/Mathematical_Thinking/designing_a_high_altitude.htm}。
ここでの参照は式の形と仮定の根拠であり、入力した$C_D$やガス量の実機に対する同定ではない。

数値APIはSciPy 1.17.0の公式マニュアル\emph{RK45}と\emph{brentq}を参照し、0.24の受入環境はSciPy 1.17.1 / NumPy 2.4.4である。
\url{https://docs.scipy.org/doc/scipy-1.17.0/reference/generated/scipy.integrate.RK45.html}、
\url{https://docs.scipy.org/doc/scipy-1.17.0/reference/generated/scipy.optimize.brentq.html}。
前者は積分・局所誤差・dense補間、後者は連続なイベント関数を挟んだ区間での根探索を担当する。
本体の試行拒否、相遷移、地形の有限延長、途中結果の保存はBalloon-JP側の契約であり、ライブラリの保証と混同しない。

GFSの配布仕様の一次入口はNOAA EMCの\url{https://www.emc.ncep.noaa.gov/emc/pages/numerical_forecast_systems/gfs.php}とNCEPの\url{https://www.nco.ncep.noaa.gov/pmb/products/gfs/nomads/}である。
JRA-3Qの原配布入口は\url{https://gdex.ucar.edu/datasets/d640000/}であり、提供DBの各値を原配布へ照合した証拠とは別である。
第\ref{th:jra-model}小節の圧力座標は、気象庁「JRA-3Q TL479フォーマット仕様」v1日本語版、第8.1節、本文pp.23--24の式（$C=1$、最上層half/2）を2026年9月30日に再読して照合した。
\url{https://www.data.jma.go.jp/jra/html/JRA-3Q/document/JRA-3Q_TL479_format_v1_ja.pdf}。
原本の係数・元日時の保持、H上の線形補間、strict支持と暫定Bは実装規約であり、気象庁による再構成精度の保証ではない。
保存形式・係数版・CLI操作は実装ノートとコマンド一覧へ分ける。

モデル・単位・補間順序・停止条件を変えるときは、対応する式、API、独立試験を同じ候補で改訂する。
PDFを独立編集せず、\path{tools/build_guides.py}でこのTeXから生成し、数式と改ページを表示確認する。
% LOGIC-END:THEORY-CORE
\end{document}
